% !TeX root = booklet-01.tex
\section{زبان ریاضی لازم برای خواندن مدار}
\label{sec:foundation-math}
فرمول یک رابطه میان کمیت‌هاست. پیش از محاسبه باید بدانیم هر حرف چه کمیتی را
نمایش می‌دهد، واحد آن چیست و رابطه تحت چه فرضی به کار می‌رود.
در این درس، حفظ شکل رابطه بدون شناخت این سه مورد می‌تواند به جواب عددی ظاهراً درست
اما از نظر فیزیکی نامعتبر منجر شود.

\subsection{نماد، زیرنویس، کسر و نسبت}
حرفی مانند \(V\) می‌تواند نماد ولتاژ باشد؛ زیرنویس آن مشخص می‌کند کدام ولتاژ
منظور است. برای مثال، \(V_a\) ولتاژ گره \(a\) و \(V_{ab}\) اختلاف ولتاژ
گره \(a\) نسبت به \(b\) است. زیرنویس، ضرب نیست: \(V_a\) به معنی \(V\times a\) نیست.
نماد \(\Delta\) نیز تغییر کمیت را نشان می‌دهد. برای نمونه، تغییر ولتاژ را
به صورت زیر تعریف می‌کنیم.
\[
 \Delta V=V_{\mathrm{final}}-V_{\mathrm{initial}}.
\]
واژه‌های زیرنویس \(\mathrm{final}\) و \(\mathrm{initial}\) به‌ترتیب حالت نهایی و اولیه‌اند.

در کسر \(a/b\)، صورت \(a\) بر مخرج \(b\) تقسیم می‌شود و باید \(b\ne0\) باشد.
نسبت دو کمیت هم‌واحد، بدون واحد است. مثلاً نسبت \(\qty{6}{\volt}\)
به \(\qty{2}{\volt}\) برابر ۳ است؛ ولی نسبت ولتاژ به جریان واحد دارد
و می‌تواند مقاومت را بیان کند. پس واحدها بخشی از استدلال‌اند.

برای دیدن اثر یک پارامتر، نخست سایر پارامترها را ثابت می‌گیریم.
در رابطه \(y=kx\)، اگر \(k\) ثابت باشد، دو برابر شدن \(x\)، مقدار \(y\) را
دو برابر می‌کند. در \(y=k/x\)، با \(k\) ثابت و \(x\ne0\)، دو برابر شدن
\(x\)، مقدار \(y\) را نصف می‌کند. این نتیجه‌ها وقتی سایر پارامترها نیز
تغییر کنند الزاماً برقرار نیستند.

\subsection{توان و نماد علمی}
توان مثبت، ضرب تکراری است. برای پایه ده و توان سه داریم
\[
 10^3=10\times10\times10=1000.
\]
توان منفی یعنی معکوس توان مثبت. برای همان پایه و توان منفی سه داریم
\[
 10^{-3}=\frac{1}{10^3}=0.001.
\]
قواعد زیر برای پایه مثبت \(a\) و توان‌های مورد استفاده در این جزوه به کار می‌روند:
\begin{equation}\label{eq:foundation-powers}
a^m a^n=a^{m+n},\qquad
\frac{a^m}{a^n}=a^{m-n},\qquad
(a^m)^n=a^{mn}.
\end{equation}
برای نمونه، قانون تفریق توان‌ها در تقسیم می‌دهد
\[
 \frac{10^{-12}}{10^{-9}}=10^{-12-(-9)}=10^{-3}.
\]
دو علامت منفی در تفریق توان‌ها حذف می‌شوند؛ نتیجه \(10^{-21}\) نیست.

در نماد علمی، عدد را به صورت \(b\times10^n\) می‌نویسیم؛ برای عدد مثبت،
معمولاً \(1\le b<10\) است. مثلاً \(0.000024=2.4\times10^{-5}\).
در جمع، باید ابتدا توان‌ها یکسان شوند:
\[
2\times10^{-3}+3\times10^{-4}
=20\times10^{-4}+3\times10^{-4}
=23\times10^{-4}=2.3\times10^{-3}.
\]
قانون جمع توان‌ها در رابطه \eqref{eq:foundation-powers} برای ضرب است؛
به جمع دو عدد تعمیم داده نمی‌شود.

\subsection{واحد و تبدیل پیشوند}
واحدهای مورد استفاده را در جدول \ref{tab:foundation-units} می‌بینیم.
نماد واحد با نماد متغیر تفاوت دارد؛ مثلاً \(C\) متغیر ظرفیت،
ولی \(\mathrm{C}\) نماد واحد بار است.
\begin{table}[ht]
\centering
\caption{کمیت‌ها و واحدهای اصلی این بخش}
\label{tab:foundation-units}
\begin{tabular}{ccc}
\toprule
کمیت & نماد متداول & واحد \lr{SI}\\
\midrule
بار & \(Q\) & \(\mathrm{C}\)\\
جریان & \(I\) & \(\mathrm{A}\)\\
ولتاژ & \(V\) & \(\mathrm{V}\)\\
مقاومت & \(R\) & \(\Omega\)\\
ظرفیت & \(C\) & \(\mathrm{F}\)\\
زمان & \(t\) & \(\mathrm{s}\)\\
انرژی & \(E\) & \(\mathrm{J}\)\\
توان & \(P\) & \(\mathrm{W}\)\\
\bottomrule
\end{tabular}
\end{table}

پیشوندهای \lr{kilo}، \lr{milli}، \lr{micro}، \lr{nano}، \lr{pico} و
\lr{femto} به‌ترتیب ضرایب \(10^3\)، \(10^{-3}\)، \(10^{-6}\)،
\(10^{-9}\)، \(10^{-12}\) و \(10^{-15}\) هستند.
پس \(\qty{2}{\kilo\ohm}=\qty{2000}{\ohm}\) و
\(\qty{20}{\femto\farad}=20\times10^{-15}\,\mathrm{F}\).

\textbf{مثال آموزشی.} برای محاسبه زمان از نسبت بار به جریان، فرض کنید
\(Q=\qty{6}{\pico\coulomb}\) و جریان ثابت \(I=\qty{3}{\milli\ampere}\) باشد:
\begin{align*}
t&=\frac{Q}{I}
=\frac{6\times10^{-12}\,\mathrm{C}}{3\times10^{-3}\,\mathrm{A}}\\
&=2\times10^{-9}\,\mathrm{s}=\qty{2}{\nano\second}.
\end{align*}
چون \(\mathrm{A}=\mathrm{C}/\mathrm{s}\)، واحد کسر برابر ثانیه می‌شود.
این کنترل واحد، مستقل از حساب عددی است و خطای استفاده از \(I/Q\) را آشکار می‌کند.

\subsection{حل معادله با حفظ تساوی}
معادله یعنی دو عبارت مقدار یکسان دارند. برای حفظ این برابری، باید عملی یکسان
روی هر دو طرف انجام دهیم. در \(V=RI\)، برای پیدا کردن جریان و با فرض \(R\ne0\)،
دو طرف را بر \(R\) تقسیم می‌کنیم تا \(I=V/R\) به دست آید.

\textbf{مثال آموزشی.} معادله \(3x+2=11\) را حل کنیم. ابتدا از دو طرف ۲ کم
می‌کنیم و سپس دو طرف را بر ۳ تقسیم می‌کنیم.
\[
 \begin{aligned}
 3x&=9,\\
 x&=3.
 \end{aligned}
\]
جای‌گذاری در معادله اولیه می‌دهد \(3\times3+2=11\).
منظور از «انتقال جمله به طرف دیگر»، همین عمل روی دو طرف است؛ علامت‌ها بی‌دلیل تغییر نمی‌کنند.

توزیع ضرب روی پرانتز یعنی \(a(b+c)=ab+ac\).
بنابراین \(2(x-3)=2x-6\)، نه \(2x-3\).
همچنین \((x-y)^2=x^2-2xy+y^2\)؛ جمله میانی را نباید حذف کرد.
برای تقسیم بر یک عبارت حاوی مجهول، باید صفر نبودن آن را بررسی کنیم؛
تقسیم معادله \(x(x-1)=0\) بر \(x\)، جواب مجاز \(x=0\) را از دست می‌دهد.

\subsection{معادله درجه‌دو و انتخاب ریشه فیزیکی}
وقتی جمله \(x^2\) داریم، ممکن است دو جواب ریاضی وجود داشته باشد.
برای \(ax^2+bx+c=0\) با \(a\ne0\)، تعریف \(D=b^2-4ac\) می‌دهد:
\begin{equation}\label{eq:foundation-quadratic}
x=\frac{-b\pm\sqrt{D}}{2a}.
\end{equation}
در اعداد حقیقی، اگر \(D<0\) باشد جواب حقیقی نداریم؛ اگر \(D=0\)،
دو ریشه بر هم منطبق‌اند و اگر \(D>0\)، دو ریشه متمایز داریم.
نماد \(\pm\) یعنی یک‌بار علامت مثبت و یک‌بار منفی را به کار ببریم.

\textbf{مثال آموزشی بدون واحد.} فرض کنید \(x^2-3x+2=0\) و مسئله شرط
\(0\le x\le1.5\) را نیز داده باشد. در اینجا \(a=1\)، \(b=-3\) و \(c=2\) است:
\[
D=(-3)^2-4(1)(2)=1,\qquad
x=\frac{3\pm1}{2}.
\]
ریشه‌ها ۱ و ۲ هستند. هر دو معادله را ارضا می‌کنند؛ فقط ریشه ۱ در بازه مجاز است.
پس جواب مسئله ۱ است، نه صرفاً ریشه مثبت؛ هر دو ریشه در این مثال مثبت‌اند.
در تحلیل ترانزیستور نیز علاوه بر معادله، باید شرط ناحیه عملکرد و بازه ولتاژها برقرار باشد.

\subsection{خواندن نمودار و معنای شیب}
در نمودار \(y=f(x)\)، محور افقی مقدار \(x\) و محور عمودی مقدار متناظر \(y\)
را نشان می‌دهد. عبارت \(f(x)\) نام یک قانون وابستگی است؛ ضرب \(f\) در \(x\) نیست.
برای دو نقطه، شیب متوسط را تعریف می‌کنیم:
\begin{equation}\label{eq:foundation-slope}
m=\frac{y_2-y_1}{x_2-x_1},\qquad x_2\ne x_1.
\end{equation}
شیب مثبت یعنی در این بازه با افزایش \(x\)، مقدار \(y\) افزایش یافته است؛
شیب منفی یعنی کاهش یافته است. واحد شیب، واحد محور عمودی تقسیم بر واحد محور افقی است.
مثلاً اگر ولتاژ در \(\qty{2}{\nano\second}\) از صفر به \(\qty{1}{\volt}\) برسد،
شیب متوسط \(\qty{0.5}{\volt\per\nano\second}\) است.

مشتق، حد همین نسبت در یک بازه بسیار کوچک است؛ نوشتن \(\mathrm{d}V/\mathrm{d}t\)
یعنی شیب لحظه‌ای نمودار ولتاژ نسبت به زمان، نه تقسیم خود \(V\) بر خود \(t\).
برای خط راست، شیب در همه نقاط ثابت و با شیب متوسط برابر است؛ برای منحنی چنین نیست.
آموزش استخراج مشتق‌ها در محل استفاده از آن‌ها ادامه خواهد یافت.

\subsection{سنجش فهم ریاضی}
\begin{enumerate}
\item \textbf{آیا \(10^{-6}+10^{-6}=10^{-12}\) است؟}
خیر؛ مجموع \(2\times10^{-6}\) است. حاصل‌ضرب آن دو برابر \(10^{-12}\) می‌شود.
\item \textbf{در \(I=V/R\)، با ثابت بودن \(V\)، سه برابر شدن مقاومت چه اثری دارد؟}
جریان به یک‌سوم می‌رسد؛ زیرا مخرج سه برابر شده است.
\item \textbf{آیا هر ریشه‌ای که معادله را ارضا کند جواب مدار است؟}
خیر؛ باید با واحدها، بازه مجاز، اتصال‌ها و فرض‌های مدل نیز سازگار باشد.
\end{enumerate}
